\documentclass[10pt]{article}
\usepackage[utf8]{inputenc}
\usepackage[T1]{fontenc}
\usepackage{amsmath}
\usepackage{amsfonts}
\usepackage{amssymb}
\usepackage[version=4]{mhchem}
\usepackage{stmaryrd}
\usepackage{bbold}

\title{PROBABILITY GU4155: Spring 2026 \\
 HOMEWORK \# 9 \\
 Professor Ioannis Karatzas \\
 Department of Mathematics, Columbia University }

\author{}
\date{}


%New command to display footnote whose markers will always be hidden
\let\svthefootnote\thefootnote
\newcommand\blfootnotetext[1]{%
  \let\thefootnote\relax\footnote{#1}%
  \addtocounter{footnote}{-1}%
  \let\thefootnote\svthefootnote%
}

%Overriding the \footnotetext command to hide the marker if its value is `0`
\let\svfootnotetext\footnotetext
\renewcommand\footnotetext[2][?]{%
  \if\relax#1\relax%
    \ifnum\value{footnote}=0\blfootnotetext{#2}\else\svfootnotetext{#2}\fi%
  \else%
    \if?#1\ifnum\value{footnote}=0\blfootnotetext{#2}\else\svfootnotetext{#2}\fi%
    \else\svfootnotetext[#1]{#2}\fi%
  \fi
}

\begin{document}
\maketitle
February 15, 2026

Consult Sections 6.1 and 6.3 of Chapter 6 in the book of John Walsh (2011).\\
Exercise \#1: If $X_{n} \rightarrow X$ and $Y_{n} \rightarrow Y$ hold in probability as $n \rightarrow \infty$, show that $X_{n} \pm Y_{n} \rightarrow X \pm Y$ and $X_{n} Y_{n} \rightarrow X Y$ also hold in probability as $n \rightarrow \infty$.

Exercise \#2: If $X_{n} \rightarrow X$ and $X_{n} \rightarrow Y$ hold in probability as $n \rightarrow \infty$, show that $\mathbb{P}(X=Y)=1$.

Exercise \#3: If $X_{n} \rightarrow c$ holds in distribution as $n \rightarrow \infty$ for some constant $c \in \mathbb{R}$, show that $X_{n} \rightarrow c$ holds also in probability as $n \rightarrow \infty$.

Exercise \#4: If $X_{n} \rightarrow 0$ holds in probability as $n \rightarrow \infty$, and $f: \mathbb{R} \rightarrow \mathbb{R}$ is bounded and uniformly continuous, show that $\mathbb{E}\left[f\left(X_{n}\right)\right] \longrightarrow f(0)$ as $n \rightarrow \infty$.

Exercise \#5: Calculate the characteristic function


\begin{equation*}
\varphi_{X}(\xi):=\mathbb{E}\left(e^{i \xi X}\right), \quad \xi \in \mathbb{R} \tag{0.1}
\end{equation*}


for a random variable $X$ with a\\
(i) Gaussian $\mathcal{N}\left(m, \sigma^{2}\right)$ distribution;\\
(i) uniform distribution on $[-a, a]$, where $a>0$.

Exercise \#6: Calculate the characteristic function $\varphi_{X}(\cdot)$ as in (0.1) above, for a random variable $X$ with\\
(i) an exponential distribution with parameter $\lambda>0$;\\
(ii) a $\operatorname{Poisson}(\lambda)$ distribution, where $\lambda>0$.

Exercise \#7: Suppose that $X, Y$ are independent random variables with common distribution $\mu$ which has mean 0 and variance 1 .

Show that $X+Y$ and $X-Y$ are independent if, and only if, the distribution $\mu$ is standard normal (Gaussian). And in this case, the random variables

$$
Z:=\frac{X+Y}{\sqrt{2}}, \quad W:=\frac{X-Y}{\sqrt{2}} \quad \text { are independent, standard normal (GAussian). }
$$

(Hint: Use without proof the TAYLOR-type expansion property $\varphi(\xi)=1-\left(\xi^{2} / 2\right)+o\left(\xi^{2}\right)$ as $\xi \rightarrow 0$, for the characteristic function $\varphi$ of $\mu$.\\
If $X+Y$ and $X-Y$ are independent, argue that this leads to the functional equation

$$
\varphi(2 \xi)=(\varphi(\xi))^{3} \varphi(-\xi), \quad \xi \in \mathbb{R}
$$

for the characteristic function $\varphi$. Now use this equation to show that $\varphi$ never vanishes; that $\rho(\xi):=\varphi(\xi) / \varphi(-\xi) \equiv 1$ holds for all $\xi \in \mathbb{R}$; and finally that all this leads to $\varphi(\xi)=e^{-\xi^{2} / 2}$.

Use without proof the fact that the Fourier-Lévy Inversion Theorem holds also for vectors of random variables: The distribution of a pair of random variables ( $Z, W$ ) is determined by its characteristic function

$$
\mathbb{R}^{2} \ni(\xi, \zeta) \mapsto \mathbb{E}\left(e^{i(\xi Z+\zeta W)}\right) \in \mathbb{C}
$$

and these two random variables are independent if, and only if, this characteristic function factors into the product of the individual characteristic functions

$$
\left.\mathbb{E}\left(e^{i(\xi Z+\zeta W)}\right)=\mathbb{E}\left(e^{i \xi Z}\right) \cdot \mathbb{E}\left(e^{i \zeta W}\right), \quad \forall(\xi, \zeta) \in \mathbb{R}^{2} .\right)
$$

Exercise \#8: For a given $\lambda>0$, let us denote by $X(\lambda)$ a random variable with a Poisson( $\lambda$ ) distribution.\\
(i) Use the Central Limit Theorem to argue that

$$
\frac{X(n)-n}{\sqrt{n}} \longrightarrow Z \sim N(0,1) \text { holds in distribution, as } n \rightarrow \infty
$$

(i) Use characteristic functions, and NOT the Central Limit Theorem, to argue that

$$
\frac{X\left(\lambda_{n}\right)-\lambda_{n}}{\sqrt{\lambda_{n}}} \longrightarrow Z \sim N(0,1) \quad \text { holds in distribution, as } n \rightarrow \infty
$$

for any sequence of positive numbers $\lambda_{n} \uparrow \infty$.\\
(Hint: At some point you may want to write a TAYLOR-type expansion of $e^{i \xi \sqrt{\lambda}}$.)\\
Exercise \#9: The Weak Law of Large Numbers, Revisited. Let $X_{1}, X_{2}, \cdots$ be independent random variables with common distribution and $\mathbb{E}\left(\left|X_{1}\right|\right)<\infty$. Denoting $S_{n}:= \sum_{j=1}^{n} X_{j}$, use characteristic functions to prove the Weak Law of Large Numbers

$$
\frac{S_{n}}{n} \rightarrow m:=\mathbb{E}\left(X_{1}\right), \quad \text { in probability, as } n \rightarrow \infty .
$$

(Hint: Use freely the first-order Taylor expansion

$$
f(x)=f(0)+x f^{\prime}(0)+x \int_{0}^{1}\left(f^{\prime}(u x)-f^{\prime}(0)\right) \mathrm{d} u, \quad x \in \mathbb{R}
$$

for the function $f(x)=e^{i x}$, in order to show

$$
\mathbb{E}\left[\exp \left\{i \frac{\xi}{n} S_{n}\right\}\right] \rightarrow e^{i \xi m} \quad \text { as } \quad n \rightarrow \infty, \quad \text { for every } \quad \xi \in \mathbb{R}
$$

Use then freely the Lévy continuity theorem for characteristic functions, as well as Exercise \#3 above, to conclude.)

Exercise \#10: The Laplace Inversion Theorem. ${ }^{1}$ Let $X: \Omega \rightarrow[0, \infty)$ be a random variable with distribution $\mu$, distribution function

$$
F(x):=\mathbb{P}(X \leq x)=\mu((-\infty, x])=\mu([0, x])=\int_{\mathbb{R}} \mathbf{1}_{[0, x]}(\theta) \mathrm{d} \mu(\theta), \quad x \in \mathbb{R}
$$

and Laplace transform

$$
\widehat{F}(\lambda):=\mathbb{E}\left(e^{-\lambda X}\right)=\int_{[0, \infty)} e^{-\lambda x} \mathrm{~d} \mu(x), \quad \lambda \in(0, \infty)
$$

Show that the Laplace transform $\widehat{F}$ determines the distribution function $F$, using the strong law of large numbers and the following steps:\\
(a) Show that the derivatives of this LAPLACE transform $\widehat{F}$ are given by

$$
\frac{\partial^{k}}{\partial \lambda^{k}} \widehat{F}(\lambda) \equiv \widehat{F}^{(k)}(\lambda)=\int_{[0, \infty)}(-1)^{k} x^{k} e^{-\lambda x} \mathrm{~d} \mu(x), \quad \lambda \in(0, \infty), k \in \mathbb{N}
$$

(In other words, justify the interchange of the operations of differentiation and integration. Once you have done this for $k=1$, it should be obvious also for $k \geq 2$.)\\
(b) Suppose that $Z_{1}(\theta), Z_{2}(\theta), \cdots$ are independent random variables, and have the same Poisson distribution with parameter $\theta=\mathbb{E}\left(Z_{k}(\theta)\right)$. Use the weak law of large numbers, to argue that

$$
\mathbb{P}\left(\frac{1}{n} \sum_{k=1}^{n} Z_{k}(\theta) \leq x\right) \longrightarrow \begin{cases}1, & \text { if } x>\theta \\ 0, & \text { if } x<\theta\end{cases}
$$

and thus

$$
\sum_{k \leq n x} e^{-n \theta} \frac{(n \theta)^{k}}{k!} \longrightarrow \begin{cases}1, & \text { if } x>\theta \\ 0, & \text { if } x<\theta\end{cases}
$$

What is the above limit, if $x=\theta$ ?\\
(c) Conclude that, for any $x \in[0, \infty)$ which is a continuity point of the distribution function $F$, we have

$$
\sum_{k \leq n x} \frac{(-1)^{k}}{k!} n^{k} \widehat{F}^{(k)}(n) \longrightarrow F(x)
$$

\footnotetext{${ }^{1}$ We stated in class, without proof, the celebrated Fourier-Lévy Inversion Theorem. This exercise deals with a "kinder, gentler" version of this result; it applies to measures supported on the positive half-line $[0, \infty)$ and thus uses Laplace, as opposed to Fourier, transforms.
}
(d) Argue that the distribution function $F$ (thus also the measure $\mu$ ) is determined uniquely from the Laplace transform $\widehat{F}$.

Exercise \#11: Do Problem 6.22, page 175 in Walsh.\\
Exercise \#12: Do Problem 6.23, page 175 in Walsh.


\end{document}
